\section{共済年金}

\begin{fillblank}[年金1 H29 問題1(2)]
  被用者年金一元化に関する次の(1)～(4)の文章について、空欄を埋めなさい。

\begin{enumerate}[(1)]
  \item  （　　A　　）年4 月に農林漁業団体職員共済組合が厚生年金に統合され、平成 27 年 10 月より国家公務
員共済組合等の共済年金が厚生年金に統一された。\\
  \item 被用者年金一元化により私立学校教職員共済制度の加入者は（　　B　　）となった。\\
  \item 厚生年金と共済年金の間の制度的な差異については、基本的には厚生年金に揃えて解消された
が、（　　C　　）はそれぞれの内容を存続している。\\
  \item 共済年金の職域部分の廃止後の新たな年金である年金払い退職給付（退職等年金給付）の給付
設計は（　　D　　）方式を採用しており、（　　E　　）、公務障害年金、公務遺族年金の３種類の給付があ
る。
\end{enumerate}
\end{fillblank}

\begin{answerbox}
  \begin{enumerate}[A: ]
    \item 平成14
    \item 第四号厚生年金被保険者
    \item 女子の支給開始年齢
    \item キャッシュバランス
    \item 退職年金
  \end{enumerate}
\end{answerbox}

\newpage